% Reconstructed text from V5 PDF page 20; non-native source
\par\smallskip\noindent\rule{\textwidth}{0.4pt}\par{\sffamily\footnotesize Original V5 page 20 -- recovered text}\par\begin{flushleft}\footnotesize\ttfamily
\par\smallskip
\par\smallskip
\noindent Diagnostics and Serving Infrastructure for MAS 20\par
\par\smallskip
\par\smallskip
\par\smallskip
\noindent optimization reduces the number of predicates by 50–65\% and rules by 50–76\%, enabling ef-\par
\noindent ficient runtime verification. The system uses short-term memory (KV-cache, discarded after\par
\noindent verification) and long-term memory (successful workflows, permanently stored), and is built on\par
\noindent the MCP framework (ICML 2025).\par
\noindent  LlamaFirewall [9] implements a three-layer defense architecture with distinct models per\par
\noindent layer. PromptGuard 2 classifies input prompts for jailbreak/injection using energy-based loss\par
\noindent with hardened tokenization, available in 86M-parameter (mDeBERTa-base) and 22M-parameter\par
\noindent (DeBERTa-xsmall) variants. AlignmentCheck evaluates each agent action against the user’s\par
\noindent stated goal using the execution trace and selected action via few-shot prompting (no fine-tuning;\par
\noindent Llama 4 Maverick recommended as backing LLM). CodeShield performs two-tiered code scan-\par
\noindent ning: Semgrep rules plus regex patterns covering 50+ CWEs across 7–8 programming lan-\par
\noindent guages. The three layers have different timing: PromptGuard is preventive (pre-processing),\par
\noindent AlignmentCheck is reactive (post-action), and CodeShield is reactive (pre-execution). The sys-\par
\noindent tem is open-source.\par
\par\smallskip
\noindent Control-theoretic recovery. ReGuard [29] applies control-theoretic principles to agent re-\par
\noindent covery: when a safety violation is detected, the system computes a corrective action using a\par
\noindent formal model of the agent’s state transition dynamics, rather than simply blocking the action.\par
\noindent This enables graceful degradation rather than hard stops, but requires a tractable state model.\par
\par\smallskip
\noindent Additional safety systems. SAIGuard [34] implements proactive defense via communication-\par
\noindent state simulation: it constructs an interaction graph over the MAS, encodes messages with\par
\noindent Qwen3-Embedding-0.6B (2-layer GNN, 1024 hidden dimensions), and estimates the impact\par
\noindent of incoming messages on local and global agent states. Risky messages are detected via recon-\par
\noindent struction deviations from benign patterns and sanitized before propagation — making SAIGuard\par
\noindent the only preventive safety system in the survey. It reduces attack success rate from 16.51\%\par
\noindent to 5.37\% (67.47\% reduction) while maintaining 89.83\% task accuracy and 98.68\% detection\par
\noindent F1. A “systemic risk amplification” finding is notable: aggregated multi-agent attack success\par
\noindent (51.78\%) far exceeds per-agent vulnerability (3.00\%).\par
\noindent  SentinelAgent [12] implements three-level anomaly detection: node-level (individual agent\par
\noindent behavior), edge-level (inter-agent communication anomalies), and path-level (multi-hop exe-\par
\noindent cution chain attacks). Path-level detection catches all 6 tested scenarios but has the highest\par
\noindent compute cost. The system uses IBM Granite Guardian 3.2 and LlamaFirewall as pre-trained\par
\noindent backends without custom training. Evaluation is qualitative only (2 case studies).\par
\noindent  AGrail [21] provides a lifelong agent guardrail framework combining rule-based constraints\par
\noindent with learned safety policies. The “lifelong” aspect enables policy updates as agents encounter\par
\noindent new scenarios without retraining. AgentDoG 1.5 [18] uses anomaly detection on execution\par
\noindent patterns to identify both safety violations and performance degradation simultaneously — the\par
\noindent only system producing both a safety verdict and a diagnostic attribution from the same analysis.\par
\noindent  ReGuard [29] applies control-theoretic recovery: when a safety violation is detected, it\par
\noindent models state transition dynamics and computes a corrective action rather than simply blocking\par
\noindent the agent. This enables graceful degradation but requires a tractable state model.\par
\par\smallskip
\noindent Diagnostic-safety overlap. Safety systems and diagnostic systems share a key requirement:\par
\noindent both need access to the agent’s execution state (tool calls, outputs, and intermediate reasoning).\par
\noindent Three of the eight safety systems (SentinelAgent, AGrail, AgentDoG) explicitly consume di-\par
\noindent agnostic traces as input to their safety decisions. This overlap suggests that serving-side state\par
\end{flushleft}
